\documentclass[11pt]{article}
\usepackage[margin=1in]{geometry}
\usepackage{amsmath,booktabs,xcolor,enumitem}
\usepackage[colorlinks=true,urlcolor=blue]{hyperref}
\newcommand{\rev}[1]{\begingroup\itshape\color{gray}#1\endgroup}
\newcommand{\resp}{\par\noindent\textbf{Response.}\ }
\newcommand{\cmt}[1]{\par\vspace{4pt}\noindent\textbf{Comment.}\ \rev{#1}\par\vspace{2pt}}
\setlength{\parskip}{4pt}
\title{\vspace{-2em}Response to Reviewers\\\large Manuscript ao-2026-07482d, ACS Omega}
\author{Krishna Harish}
\date{}
\begin{document}
\maketitle
\vspace{-2em}

\noindent\textbf{Title:} Objective-Dependent Active Learning and Calibrated
Uncertainty for Sample-Efficient Discovery of Rare-Earth Extractants: A Benchmark on
1{,}202 Experimental Distribution Coefficients with Prospective Validation.

We thank the reviewer for a thorough and constructive report. Every comment is answered
below \emph{in the order it was raised}, quoting the comment and stating our response
and the specific change made. All methodological requests are addressed with
\textbf{new experiments on the real data}; no claim is asserted without a computed
result. The revision now reports five findings, nine figures, four tables, and 17 added
references, all regenerated by released scripts. Section, figure, and table numbers
refer to the revised manuscript.

\section*{Reviewer 1 --- overall assessment}

\cmt{Machine learning assisted rare earth solvent extraction is studied in the research
by benchmarking active learning strategies on experimental logD data and by explicitly
separating two practical objectives: building an accurate predictive model and
discovering high performing extractants.}
\resp We thank the reviewer for the accurate summary. The prediction-versus-discovery
separation remains the organising theme of the paper and is now tested more
rigorously (leakage-controlled splits, a ligand-batch acquisition unit, and
definition-resolved recall).

\cmt{The study notes that active learning does not automatically improve prediction,
exploitation biased acquisition can recover top candidates while degrading global model
quality, and uncertainty calibration is essential for trustworthy reporting even when
it does not improve acquisition ranking. This objective dependent framing is a notable
aspect of the paper and is useful for researchers planning expensive extraction
experiments.}
\resp We appreciate this. These three points are retained and strengthened; the
calibration point is now also shown to hold under a realistic ligand-disjoint split,
with the added caveat that non-monotone calibration can change acquisition (see the
calibration comment below).

\cmt{The use of a random forest baseline in the study is also worth note. Showing that a
relatively simple ensemble model can outperform a previously reported deep network on
the same feature representation is important, particularly for tabular chemical datasets
with moderate sample size.}
\resp Thank you. The random-forest-versus-deep-network comparison is retained on the
same published row-level split (an apples-to-apples model comparison). We now add,
immediately after it, that this in-distribution number is optimistic relative to
leakage-controlled splits and is not a deployability estimate (next comment).

\cmt{The multi seed active learning curves, comparison among random, greedy,
uncertainty, UCB, EI, and hybrid acquisition, and replication on an independent
lipophilicity benchmark all support the argument that prediction oriented and discovery
oriented acquisition can diverge. The prospective test on newly synthesised ligands is
sound, since external validation is still uncommon in extractant property modelling.}
\resp We thank the reviewer. All of these are retained; the six-strategy comparison is
additionally reported under the ligand-batch acquisition unit and under three
definitions of a ``top'' system (below).

\section*{Reviewer 1 --- methodological points}

\cmt{The experimental rows combine ligand identity, lanthanide identity, solvent
conditions, and measured logD. If train/test splits are made at the row level, closely
related measurements from the same ligand or condition family can appear in both
training and testing, potentially inflating validation performance. A more stringent
split by ligand scaffold, publication source, or ligand family would better represent
discovery of genuinely new extractants. The drop from in-domain R2 = 0.94 to
prospective R2 = 0.69 already suggests a meaningful chemotype shift effect, so scaffold
level validation would provide a more realistic estimate of deployability.}
\resp Addressed, and now a central result (new \S``Realistic generalisation under
leakage-controlled splits,'' Fig.~1, Table~1). Although the released matrix has no
SMILES, the grouping structure is recoverable: identical ECFP bit-vectors identify the
93 ligands, the metal-descriptor block the 14 lanthanides, and the citation index the
45 sources. Under leakage-controlled splits:

\begin{center}\small
\begin{tabular}{lc}
\toprule
Split & RF $R^2$ \\ \midrule
Random (row-level) & $0.88\pm0.03$ \\
Ligand-disjoint & $0.20\pm0.31$ \\
Source-disjoint & $-0.03\pm0.80$ \\
Ligand-family (ECFP-Tanimoto, scaffold proxy) & $-0.24\pm0.50$ \\
Prospective (4 new ligands) & $0.70$ \\
\bottomrule
\end{tabular}
\end{center}

\noindent The row-level number is strongly optimistic; grouped splits collapse it to at
or below the mean. We now state explicitly that in-distribution $R^2$ is not a
deployability estimate, and that the prospective $R^2=0.70$ sits between the extremes
because the four new ligands are close analogues of training chemotypes. We cite
Sheridan 2013 and Guo et al.\ 2024/2025. Because SMILES are unavailable, we are
transparent (Limitations) that the family split is an ECFP-Tanimoto proxy for a
scaffold split.

\cmt{The active learning formulation should also clarify the true experimental decision
unit. In laboratory extractant discovery, selecting one ligand can imply measuring
multiple lanthanides, acidities, diluents, or concentrations, rather than selecting
isolated ligand--metal--condition rows. Ranking and acquiring individual rows can
therefore overstate operational flexibility. A ligand batch acquisition simulation,
where acquisition adds all measurements associated with a selected ligand or a minimal
experimental panel, would make the benchmark more chemically realistic.}
\resp Added (new \S``The dichotomy survives a realistic ligand-batch unit,'' Fig.~4). We
re-ran the benchmark with a ligand-batch unit: selecting a ligand reveals all of its
measurements (median 14, across lanthanides/acidities/diluents/temperatures), and the
test set is a disjoint set of held-out ligands. The design-side dichotomy is clear:
exploitation gives the best top-10 ligand recall (1.00) while random/exploration lags
(0.83). On prediction, all $R^2$ are low under this hard ligand-disjoint test, with
random best (0.30) by a small margin over exploitation (0.29). We report this honestly:
the objective-dependent ordering is preserved under the operational decision unit,
though the prediction-side gap narrows once train--test ligand overlap is forbidden.

\cmt{The `top 10 recall' metric should specify whether `top' systems are unique ligands,
ligand--metal pairs, or condition specific measurements; discovery value differs
substantially across these definitions.}
\resp Added (new \S``What counts as a `top' system,'' Fig.~5, Table~2). Recall is now
reported under all three definitions, and the ranking of strategies inverts:
exploitation dominates at the row level (greedy 1.00 vs random 0.45), but at the ligand
level broad sampling already recovers most top ligands (random 0.95 vs greedy 0.82; EI
best at 0.97), because each ligand appears under many conditions. We caution that an
unstated ``top-$k$ recall'' can reverse the practical conclusion.

\cmt{The uncertainty calibration section is important, but the conformal procedure needs
more precise reporting. Calibration splits should be fully separated from both training
and final testing within each active learning budget, and any rescaling must avoid
leakage from the held out pool. Reporting coverage by chemical subgroup, logD range,
lanthanide identity, and prospective ligands would be more informative than global
coverage alone, because well calibrated average intervals may still fail in the high
logD region most relevant to discovery. For active learning, rank preservation after
monotonic calibration explains the unchanged acquisition performance, but this point
should be stated with the caveat that non monotonic or locally adaptive calibration
could alter rankings.}
\resp Fully reworked (new \S, Fig.~6, Table~3), addressing each part in turn.
(i)~\emph{Separation of splits}: within the active-learning loop the calibration set is
now drawn from the already-revealed pool at each budget and is disjoint from the fixed
held-out test set (stated in the SI Methods); the subgroup analysis additionally uses a
ligand-disjoint train/calibration/test partition, so calibration, training, and testing
share no ligand, and the rescaling cannot leak from the held-out pool.
(ii)~\emph{Subgroup coverage}: we report coverage of nominal-80\% intervals by logD
region, lanthanide class, and the prospective set, not only globally: raw intervals
under-cover (0.72), a single global conformal factor restores average coverage (0.90)
but leaves the high-logD region uneven, and Mondrian (group-conditional) conformal
(0.87) evens the subgroups; per-lanthanide-class coverage is 0.89--0.91 and prospective
0.88. (iii)~\emph{High-logD failure}: this is exactly where global conformal is weakest
and where Mondrian helps (Table~3). (iv)~\emph{Non-monotone caveat}: we now state that
a rank-preserving global factor cannot change acquisition, \emph{but} Mondrian reorders
the uncertainty ranking (Spearman $\rho=0.85<1$), so locally adaptive calibration
\emph{can} change acquisition. We retain the in-distribution result (calibration error
$0.103\!\to\!0.013$) and frame the two settings as complementary.

\cmt{The study also needs broader baseline comparisons. Random forest is a strong model,
but Gaussian process regression, quantile random forests, gradient boosting, Bayesian
neural networks, and conformalized residual models could provide more insight into
uncertainty ranking quality. Reporting top-k enrichment, regret, hit rate above a
practical logD threshold, and Pareto performance for selectivity relevant objectives
would strengthen the connection to separations chemistry, where a high distribution
coefficient alone is not always sufficient.}
\resp Addressed in three parts. (i)~\emph{Models}: we added a five-estimator comparison
(new \S, Fig.~7, Table~4): random forest, quantile regression forest, Gaussian process
(PCA-reduced), gradient boosting, and an MC-dropout Bayesian neural network, each with a
split-conformal (conformalized-residual) wrapper, on a ligand-disjoint split, scored on
predictive accuracy, conformal coverage, and uncertainty-ranking quality (Spearman of
$\sigma$ vs.\ error; selective-prediction RMSE-AUC). No estimator uniformly dominates,
with high fold variance; gradient boosting is most consistent (best mean ranking quality
$\rho=0.28$) and the Gaussian process is weakest. (ii)~\emph{Discovery metrics}: we
added the enrichment factor, simple regret, and hit-rate above a practical threshold
($\log D\ge1$, i.e.\ $D\ge10$) to Table~2, plus the $R^2$-vs-recall Pareto view
(Figs.~3--5). (iii)~\emph{Selectivity}: on the point that ``a high distribution
coefficient alone is not always sufficient,'' we added a selectivity analysis (new \S,
Fig.~8): per ligand we compute the separation factor as the spread of $\log D$ across
the 14 lanthanides (up to 6.1 log units over 63 ligands). Ranking by extraction strength
(max $\log D$) is not the same as ranking by selectivity (Spearman 0.63); only 4 of the
top-10 strongest extractants are also among the top-10 most selective, and we give the
strength--selectivity Pareto front. A full selectivity-aware acquisition loop is noted
as a natural extension.

\section*{Reviewer 1 --- remaining points}

\cmt{Major revision needed before the paper can be accepted for publication.}
\resp We have carried out the major revision described above.

\cmt{Language use is mainly correct.}
\resp Thank you; we have additionally proof-read the revised text.

\cmt{More Figures are needed to enhance the technical quality of the paper.}
\resp We added five new figures (leakage-controlled splits, ligand-batch acquisition,
recall definitions, subgroup calibration, model/UQ comparison) plus the selectivity
figure, for nine figures in total, all auto-generated from the results files.

\cmt{Include tables in the paper to support reading flow and technical quality.}
\resp We added four tables (splits, discovery metrics, subgroup calibration, model/UQ
comparison), each generated from the released results.

\cmt{More literature study is needed to position the paper in the broader body of
literature.}
\resp The introduction and discussion now cite 17 added, verified references
(2013--2025): pool-based/batch AL (Graff 2021, Rohr 2020, Reker 2020, Bayley 2024,
Bailey 2023), explore/exploit (De Ath 2021), realistic splitting (Sheridan 2013, Guo
2024/2025), group-conditional conformal (Sun 2017), early-recognition metrics
(Truchon \& Bayly 2007, Ash 2022), UQ-estimator comparisons (Hirschfeld 2020, Scalia
2020, Soleimany 2021, Busk 2022, Gal 2016), and current rare-earth $\log D$ ML
(Udawattha \& Alam 2025, Chaube 2020, Edaugal 2025).

\vspace{6pt}
\noindent We hope the reviewer finds that every point has been addressed with concrete,
reproducible evidence and that the revised manuscript is substantially strengthened.

\end{document}
